\subsection{Change of rings}

PROPOSITION 5.--- Let \(\rho: A \to B\) a homomorphism of noetherian rings, making B a flat A-module. One assumes that for every maximal ideal \(\mathfrak{n}\) of B, the ring \(\kappa(\rho^{-1}(\mathfrak{n})) \otimes_{A} B\) is a Gorenstein ring. Let \(\Omega\) be a dualizing A-module; the B-module \(\Omega_{(B)}\) is dualizing.

Let \(\mathfrak{n}\) a maximal ideal of B, and \(\mathfrak{p}\) its inverse image in A. The \(A_{\mathfrak{p}}\) -module \(B_{\mathfrak{n}}\) is flat, the \(A_{\mathfrak{p}}\) -module \(\Omega_{\mathfrak{p}}\) is dualizing, \(\Omega_{(B)} \otimes_B B_{\mathfrak{n}}\) is identified with \(\Omega_{\mathfrak{p}} \otimes_{A_{\mathfrak{p}}} B_{\mathfrak{n}}\) and \(\kappa_{\Lambda_{\mathfrak{p}}} \otimes_{\Lambda_{\mathfrak{p}}} B_{\mathfrak{n}}\) , which is identified with a ring of fractions of \(\kappa(\mathfrak{p}) \otimes_A B\) , is a Gorenstein ring. It therefore suffices to prove the proposition when \(\rho\) is a local homomorphism of local rings, which we will assume from now on.

Let us first treat the case where the rings A and B are Artinian. Set \(C = B / \mathfrak{m}_{\Lambda}B\) Since B is flat over A, the B-module \(\mathrm{Hom}_{\mathrm{B}}(\mathrm{C},\Omega_{(\mathrm{B})})\) is isomorphic to \(\mathrm{Hom}_{\Lambda}(\kappa_{\mathrm{A}},\Omega)\otimes_{\mathrm{A}}\mathrm{B}\) (I, § 2, no. 10, prop. 11), hence to \(\mathrm{Hom}_{\mathrm{A}}(\kappa_{\mathrm{A}},\Omega)\otimes_{\kappa_{\mathrm{A}}}\mathrm{C}\) From this we deduce a sequence of isomorphisms

\[
\begin{array}{c} \operatorname{Hom} _ {\mathrm{B}} (\kappa_ {\mathrm{B}}, \Omega_ {(\mathrm{B})}) \longrightarrow \operatorname{Hom} _ {\mathrm{C}} (\kappa_ {\mathrm{C}}, \operatorname{Hom} _ {\mathrm{B}} (\mathrm{C}, \Omega_ {(\mathrm{B})})) \longrightarrow \operatorname{Hom} _ {\mathrm{C}} (\kappa_ {\mathrm{C}}, \operatorname{Hom} _ {\mathrm{A}} (\kappa_ {\mathrm{A}}, \Omega) \otimes_ {\kappa_ {\mathrm{A}}} \mathrm{C}) \\ \qquad \qquad \qquad \qquad \qquad \qquad \qquad \qquad \qquad \qquad \longrightarrow \operatorname{Hom} _ {\mathrm{A}} (\kappa_ {\mathrm{A}}, \Omega) \otimes_ {\kappa_ {\mathrm{A}}} \operatorname{Hom} _ {\mathrm{C}} (\kappa_ {\mathrm{C}}, \mathrm{C}). \end{array}
\]

The \(\kappa_{\mathrm{A}}\) -vector space \(\mathrm{Hom}_{\mathrm{A}}(\kappa_{\mathrm{A}}, \Omega)\) is of dimension 1 since \(\Omega\) is dualizing, and the same holds for the \(\kappa_{\mathrm{C}}\) -vector space \(\mathrm{Hom}_{\mathrm{C}}(\kappa_{\mathrm{C}}, \mathrm{C})\) since C is a Gorenstein ring; consequently, \(\kappa_{\mathrm{B}}\) -vector space \(\mathrm{Hom}_{\mathrm{B}}(\kappa_{\mathrm{B}}, \Omega_{(\mathrm{B})})\) has dimension 1.

Let M be a B-module of finite length; let us prove by induction on \(\operatorname{long}_{\mathrm{B}}(\mathrm{M})\) that we have \(\operatorname{long}_{\mathrm{B}}(\operatorname{Hom}_{\mathrm{B}}(\mathrm{M},\Omega_{(\mathrm{B})})) \leqslant \operatorname{long}_{\mathrm{B}}(\mathrm{M})\) The assertion is clear if M = 0, and it follows from what precedes if \(M = \kappa_{B}\) . Suppose \(\operatorname{long}_{\mathrm{B}}(\mathrm{M}) \geqslant 2\) There exists an exact sequence of B-modules

\[
0 \rightarrow \mathrm{M} ^ {\prime} \rightarrow \mathrm{M} \rightarrow \kappa_ {\mathrm{B}} \rightarrow 0
\]

with \(\mathrm{long}_{\mathrm{B}}(\mathrm{M}^{\prime}) < \mathrm{long}_{\mathrm{B}}(\mathrm{M})\) One deduces an exact sequence

\[
0 \to \operatorname{Hom} _ {\mathrm{B}} (\kappa_ {\mathrm{B}}, \Omega_ {(\mathrm{B})}) \to \operatorname{Hom} _ {\mathrm{B}} (\mathrm{M}, \Omega_ {(\mathrm{B})}) \to \operatorname{Hom} _ {\mathrm{B}} (\mathrm{M} ^ {\prime}, \Omega_ {(\mathrm{B})}),
\]

and one concludes by applying the induction hypothesis to M'.

Let N be the kernel of the canonical surjection from \(\kappa_{A} \otimes_{A} B\) on \(\kappa_{B}\) . Set \(m = \text{long}_{\text{B}}(\kappa_{\text{A}} \otimes_{\text{A}} \text{B})\) ; one has \(\text{long}_{\text{B}}(\text{N}) = m - 1\) . Consider the exact sequence of B-modules

\[
\begin{array}{c} 0 \to \operatorname{Hom} _ {\mathrm{B}} (\kappa_ {\mathrm{B}}, \Omega_ {(\mathrm{B})}) \longrightarrow \operatorname{Hom} _ {\mathrm{B}} (\kappa_ {\mathrm{A}} \otimes_ {\mathrm{A}} \mathrm{B}, \Omega_ {(\mathrm{B})}) \longrightarrow \operatorname{Hom} _ {\mathrm{B}} (\mathrm{N}, \Omega_ {(\mathrm{B})}) \\ \longrightarrow \operatorname{Ext} _ {\mathrm{B}} ^ {1} (\kappa_ {\mathrm{B}}, \Omega_ {(\mathrm{B})}) \longrightarrow \operatorname{Ext} _ {\mathrm{B}} ^ {1} (\kappa_ {\mathrm{A}} \otimes_ {\mathrm{A}} \mathrm{B}, \Omega_ {(\mathrm{B})}). \end{array}
\]

The B-modules \(\mathrm{Hom}_{\mathbb{B}}(\kappa_{\Lambda}\otimes_{\Lambda}\mathrm{B},\Omega_{(\mathbb{B})})\) and \(\mathrm{Ext}_{\mathbb{B}}^{1}(\kappa_{\Lambda}\otimes_{\Lambda}\mathrm{B},\Omega_{(\mathbb{B})})\) are respectively isomorphic to \(\mathrm{Hom}_{\Lambda}(\kappa_{\Lambda},\Omega)\otimes_{\Lambda}\mathrm{B}\) and \(\mathrm{Ext}_{\mathbb{A}}^{1}(\kappa_{\Lambda},\Omega)\otimes_{\Lambda}\mathrm{B}\) , that is, to \(\kappa_{\Lambda}\otimes_{\Lambda}\mathrm{B}\) and to 0. The lengths of the B-modules \(\mathrm{Hom}_{\mathbb{B}}(\kappa_{\mathbb{B}},\Omega_{(\mathbb{B})})\) and \(\mathrm{Hom}_{\mathbb{B}}(\kappa_{\mathbb{A}}\otimes_{\mathbb{A}}\mathrm{B},\Omega_{(\mathbb{B})})\) are 1 and \(m\) and that of \(\mathrm{Hom}_{\mathbb{B}}(\mathrm{N},\Omega_{(\mathbb{B})}))\) is \(\leqslant m - 1\) from which it follows that the B-module \(\mathrm{Ext}_{\mathbb{B}}^{1}(\kappa_{\mathbb{B}},\Omega_{(\mathbb{B})})\) is zero. By prop. 6 of § 3, no. 3, the B-module \(\Omega_{(\mathbb{B})}\) is injective; consequently it is a dualizing module (no. 1, example 1).

Moving to the general case, we set \(C = \kappa_A \otimes_A B\) ; by hypothesis, this is a Gorenstein ring, hence a Macaulay ring ( \(\S\) By prop. 1 of no. 1, A is a Macaulay ring, and the A-module \(\Omega\) is Macaulay. Consequently B is a Macaulay ring, and the B-module \(\Omega_{(B)}\) is Macaulay \(\S\) 2, no. 7, cor. 1 of prop. 9). Set \(r = \dim(A)\) , \(s = \dim(C)\) There exists a sequence ( \(x_1, \ldots, x_r\) ) of elements of \(\mathfrak{m}_A\) regular for the A-modules A and \(\Omega\) and a sequence ( \(y_1, \ldots, y_s\) ) of elements of \(\mathfrak{m}_B\) regular for the B-module C; denote \(x\) denote the ideal of A and \(y\) the ideal of B that they generate respectively. The sequence ( \(y_1, \ldots, y_s, \rho(x_1), \ldots, \rho(x_r)\) ) is regular for the B-modules B and \(\Omega_{(B)}\) ( \(\S\) 1, no. 6, prop. 11), and the A-module B/ \(y\) is flat (loc. cit., prop. 10). Set \(A' = A/x\) , \(B' = B/(xB + y)\) and let \(\rho': A' \to B'\) the homomorphism deduced from \(\rho\) by passing to quotients. The rings \(A'\) and \(B'\) are Artinian, the A'-module \(B'\) is flat, the ring \(\kappa_{A'} \otimes_{A'} B'\) , which is identified with C/ \(y\) is a Gorenstein ring ( \(\S\) 3, no. 7, example 2) and the A'-module \(\Omega_{(A')}\) is dualizing (no. 2, prop. 4). By the first part of the proof, the B'-module \(\Omega_{(B')}\) is dualizing. It then follows from loc. cit. that the B-module \(\Omega_{(B)}\) is dualizing.

COROLLARY.---Let A be a Noetherian ring admitting a dualizing module. \(\Omega\) B be a polynomial algebra over A in a finite number of indeterminates. The B-module \(\Omega_{(B)}\) is dualizing.

Indeed, for every prime ideal p of A, the ring \(\kappa(\mathfrak{p})\otimes_{\mathrm{A}}\mathrm{A}[\mathbf{X}]\) is identified with \(\kappa(\mathfrak{p})[\mathbf{X}]\) , which is regular, hence Gorenstein.

PROPOSITION 6. — Let \(\Lambda\) be a Noetherian local ring and \(\Omega\) a dualizing A-module. Let B be a finite A-algebra; assume that the A-module B is Macaulay. The B-module \(\mathrm{Ext}_{\mathrm{A}}^{i}(\mathrm{B},\Omega)\) is zero for \(i\neq \dim (\Lambda) - \dim (\mathrm{B})\) and dualizing for \(i = \dim (\Lambda) - \dim (\mathrm{B})\) .

We have \(\dim(B) = \dim_{A}(B) \leqslant \dim(\Lambda)\) (VIII, § 2, no. 3, th. 1 c)); set \(c = \dim(A) - \dim(B)\) . One has \(\operatorname{Ext}_{A}^{i}(B, \Omega) = 0\) for \(i \neq c\) since the A-module B is Macaulay (no. 1, cor. of prop. 3). Let us prove that the B-module \(\operatorname{Ext}_{\Lambda}^{c}(B, \Omega)\) is dualizing.

Suppose first \(\dim(B) = 0\) . The spectrum X of B is finite and consists of maximal ideals (IV, § 2, no. 5, prop. 9); the canonical map \(B \to \prod_{n \in X} B_n\) is an isomorphism (loc. cit., cor. 1). The B-module \(\Omega' = \text{Ext}_A^c(B, \Omega)\) is therefore a direct sum of the modules \(\text{Ext}_A^c(B_n, \Omega)\) ; as \(\text{Ext}_A^c(B_n, \Omega)\) has support in \(\{\mathfrak{n}\}\) , it is identified with \(\Omega'_n\) . One has \(\dim(B_n) = 0\) for every \(\mathfrak{n}\) ; to prove that the B-module \(\Omega'\) is dualizing, it therefore suffices to prove that this is so for the \(B_n\) -module \(\text{Ext}_A^c(B_n, \Omega)\) for every \(\mathfrak{n} \in X\) , which reduces us to the case where the ring B is local. In this case, according to example 6 of § 8, no. 5, the B-module \(\text{Ext}_A^c(B, \Omega)\) is isomorphic to \(\text{Hom}_A(B, I)\) , where I is a Matlis A-module; it is consequently a Matlis B-module (§ 8, no. 6, cor. of prop. 5), hence a dualizing B-module (no. 1, example 1).

Now suppose \(\dim(B) > 0\) and reason by induction on \(\dim(B)\) . One has \(\operatorname{prof}_{A}(B) = \dim_{A}(B) = \dim(B)\) , whence \(\operatorname{prof}_{A}(B) > 0\) ; on the other hand we have \(\operatorname{prof}(A) = \dim(A) > 0\) (no. 1, prop. 1), and consequently \(\operatorname{prof}_{A}(A \oplus B) > 0\) . There therefore exists an element \(x\) of \(\mathfrak{m}_{\Lambda}\) such that the homotheties \(x_{\Lambda}\) and \(x_{B}\) are injective.

Consider the exact sequence of extension modules associated with the exact sequence \(0 \to \mathrm{B} \xrightarrow{x_{\mathrm{B}}} \mathrm{B} \longrightarrow \mathrm{B}/x\mathrm{B} \to 0\) and with the A-module \(\Omega\) . The A-module B/xB is Macaulay ( \(\S 2, n^{\circ} 1\) , example 3), of dimension \(\dim(\mathrm{B}) - 1\) (VIII, \(\S 3, n^{\circ} 2\) ; we therefore have \(\operatorname{Ext}_{\mathrm{A}}^{i}(\mathrm{B}/x\mathrm{B}, \Omega) = 0\) for \(i \neq c + 1\) ( \(n^{\circ} 1\) (corollary of Prop. 3). Since we have \(\operatorname{Ext}_{\mathrm{A}}^{i}(\mathrm{B}, \Omega) = 0\) for \(i \neq c\) , one obtains an exact sequence of B-modules

\[
0 \rightarrow \operatorname{Ext} _ {\mathrm{A}} ^ {c} (\mathrm{B}, \Omega) \xrightarrow {x} \operatorname{Ext} _ {\mathrm{A}} ^ {c} (\mathrm{B}, \Omega) \longrightarrow \operatorname{Ext} _ {\mathrm{A}} ^ {c + 1} (\mathrm{B} / x \mathrm{B}, \Omega) \rightarrow 0.
\]

By the induction hypothesis, the B/xB-module \(\operatorname{Ext}_{\Lambda}^{c+1}(\mathrm{B}/x\mathrm{B},\Omega)\) is dualizing. Since the A-algebra B is finite, the image of \(m_{A}\) in B is contained in the radical of B (V, § 2, no. 1, prop. 1); by prop. 4 of no. 2, the B-module \(\operatorname{Ext}_{\mathrm{A}}^{c}(\mathrm{B},\Omega)\) is dualizing.

COROLLARY 1.--- Let A be a Noetherian ring, Ω a dualizing A-module, and B a finite A-algebra; suppose that the A-module B is Macaulay. The B-module Ext \(_{A}\) (B, Ω) is dualizing.

Let \(\Omega'\) the B-module \(\mathrm{Ext}_{\mathsf{A}}(\mathsf{B},\Omega)\) . Let \(\mathfrak{n}\) a maximal ideal of B; its inverse image in A is a maximal ideal \(\mathfrak{m}\) (V, § 2, no. 1, prop. 1). The \(\Lambda_{\mathfrak{m}}\) -algebra \(\mathrm{B}_{\mathfrak{m}} = \mathrm{A}_{\mathfrak{m}} \otimes_{\mathsf{A}} \mathrm{B}\) is finite, and it is a \(\mathrm{A}_{\mathfrak{m}}\) -module macaulayen; by the proposition, the \(\mathrm{B}_{\mathfrak{m}}\) -module \(\Omega_{\mathfrak{m}}'\) , which is identified with \(\mathrm{Ext}_{\mathrm{A}_{\mathfrak{m}}}(B_{\mathfrak{m}},\Omega_{\mathfrak{m}})\) (§ 3, no. 2, prop. 2) is dualizing. Since \(\mathrm{B}_{\mathfrak{n}}\) is a ring of fractions of \(\mathrm{B}_{\mathfrak{m}}\) , the \(\mathrm{B}_{\mathfrak{n}}\) -module \(\Omega_{\mathfrak{n}}'\) is dualizing, whence the corollary.

Remark. Keep the hypotheses of Cor. 1 and suppose further that the canonical homomorphism \(\rho : A \to B\) be injective. We then have \(\dim(\Lambda_{\mathfrak{m}}) = \dim(B_{\mathfrak{m}})\) for every maximal ideal \(\mathfrak{m}\) of \(A\) By prop. 6 and cor. 1, \(\operatorname{Ext}_{A}^{i}(B, \Omega)\) is zero for \(i \neq 0\) and the B-module \(\operatorname{Hom}_{A}(B, \Omega)\) is dualizing.

COROLLARY 2.--- If a Noetherian ring A possesses a dualizing module, then every finitely generated A-algebra which is a Macaulay ring possesses a dualizing module.

This follows from cor. 1 and from the corollary of prop. 5.

COROLLARY 3.--- Every presentable Macaulay ring (in particular, every complete local Macaulay ring) possesses a dualizing module.

Indeed, let R be a regular ring and A a Macaulay ring quotient of R. The R-module A is Macaulay (§ 2, no. 5, example 5), and R possesses a dualizing module (no. 1, example 2); hence so does Λ by cor. 1. Moreover, it has already been observed that a complete local Noetherian ring is presentable (§ 4, no. 4, prop. 6, c)).

More generally, every Macaulay ring that is a quotient of a Gorenstein ring has a dualizing module. Conversely, one can show that a local Macaulay ring that has a dualizing module is a quotient of a local Gorenstein ring (exerc. 1).

